\documentclass[12pt]{article}
\usepackage[a4paper,margin=2.5cm]{geometry}
\usepackage{amsmath,amssymb}
\usepackage{fontspec}

\newfontfamily\cjk{Noto Sans CJK HK}
\newcommand{\zh}[1]{{\cjk #1}}

\setlength{\parindent}{0pt}

\begin{document}

{\large\textbf{\zh{Step 1}}}\\[0.5em]
\zh{模型取樣範圍,半徑 } $[40\,\mu\text{m},\ 500\,\mu\text{m}]$\\[0.3em]
\zh{這半徑區間要幾乎連續,接著請從這些取樣範圍提取各別的}
$\widehat{\mathbf{K}}_I^{\,\mathrm{FEM}}$,$\widehat{g}_B$,$\widehat{\ell}$\,。

\vspace{1.5em}

{\large\textbf{\zh{Step 2}}}\\[0.5em]
\zh{出一張疊圖,橫軸(取樣半徑),縱軸(誤差),而誤差定義,}
\zh{以} $\widehat{\mathbf{K}}_I^{\,\mathrm{FEM}}$\,:

\vspace{1em}

\[
\left\{
\begin{aligned}
\widehat{\mathbf{K}}_I^{\,\mathrm{FEM}} \;&\Rightarrow\; \text{Frobenius }\text{\cjk 範數}\\[1.5em]
\widehat{g}_B \;&\Rightarrow\; \frac{\widehat{g}_{Bi}-\widehat{g}_{Bj}}{\widehat{g}_{Bj}}\\[1.5em]
\widehat{\ell} \;&\Rightarrow\; \frac{\widehat{\ell}_{i}-\widehat{\ell}_{j}}{\widehat{\ell}_{j}}
\end{aligned}
\right.
\]

\vspace{1em}

\zh{其中}
$\left\{
\begin{aligned}
&i:\ \text{\cjk 後一步取樣範圍(從 } 40\,\mu\text{m}\ \text{\cjk 開始取)}\\
&j:\ \text{\cjk 前一步取樣範圍}
\end{aligned}
\right.$

\end{document}
